\documentclass[10pt]{article}
\usepackage[margin=0.7in]{geometry}
\usepackage{amsmath,amssymb,mathtools,hyperref}
\hypersetup{colorlinks=true,urlcolor=blue,linkcolor=blue}
\setlength{\parindent}{0pt}
\setlength{\parskip}{3pt}
\begin{document}
\begin{center}
{\Large\bfseries Problem 00000001581: rational and finite-field points}\\[5pt]
lizaixi01 \quad AI-led research and formalization \quad 23 October 2026
\end{center}
\setlength{\abovedisplayskip}{6pt}
\setlength{\belowdisplayskip}{6pt}

\textbf{Result and scope.} The source explicitly says ``Q-points'' (also $\mathbb Q$ in Chinese), but gives the finite expression $p^2+ap+1$, $a\in\{-1,0,1\}$. No PVI parameters are specified. We treat both the literal rational reading and the finite-field reading used in the earlier review. For the standard zero-local-trace PVI character surface,
\[
 S:\quad x^2+y^2+z^2+xyz=4,
\]
we prove that $S(\mathbb Q)$ is infinite, even on its Jacobian-regular locus, and that $\#S(\mathbb F_p)=p^2+1$ for \emph{every} prime $p$. Thus the literal finite rational-point claim is false, while the intended finite-field pattern is proved for this actual PVI member, with $a=0$. We do not claim a formula for every other PVI parameter.

\section*{1. Identifying the actual PVI member}
Cantat--Loray \cite{CL}, printed pp.~2928--2929, use the four-punctured-sphere presentation
\[
 \langle\alpha,\beta,\gamma,\delta\mid\alpha\beta\gamma\delta=1\rangle.
\]
For an $\mathrm{SL}_2$ representation with local traces $a,b,c,d$, set
$x=\operatorname{tr}(\alpha\beta)$, $y=\operatorname{tr}(\beta\gamma)$, and $z=\operatorname{tr}(\gamma\alpha)$. Their equations (1.3)--(1.5) identify the character surface as
\[
 x^2+y^2+z^2+xyz=Ax+By+Cz+D,
\]
where $A=ab+cd$, $B=bc+ad$, $C=ac+bd$, and
$D=4-a^2-b^2-c^2-d^2-abcd$. Local traces all zero give exactly $S$. Equation (1.12), p.~2932, gives local traces $2\cos(\pi\theta_i)$, so $\theta_i=1/2$ for all four punctures is a legitimate PVI choice.

We adopt this published standard character-surface model. Lean independently constructs the displayed presented group, actual determinant-one matrices, their homomorphism, and the actual trace coordinates below. It does not reprove the complete topological fundamental-group identification, GIT generation theorem, or Riemann--Hilbert correspondence.

\section*{2. Infinitely many genuine rational trace points}
For $n\geq0$, put $u=n+2$ and $s=u+u^{-1}$. Use
\[
 A=\begin{pmatrix}0&-1\\1&0\end{pmatrix},\quad B=-A,\quad
 C=\begin{pmatrix}0&u\\-u^{-1}&0\end{pmatrix},\quad D=-C.
\]
All four matrices lie in $\mathrm{SL}_2(\mathbb Q)$ and have trace zero. Since $AB=CD=I$, their product is $I$; consequently they define an actual representation of the presented group. Direct multiplication gives
\[
 (\operatorname{tr}(AB),\operatorname{tr}(BC),\operatorname{tr}(CA))=(2,-s,s)\in S(\mathbb Q).
\]
These are distinct: $u>1$, and for $v>u>1$,
\[
 (v+v^{-1})-(u+u^{-1})=\frac{(v-u)(uv-1)}{uv}>0.
\]
Also $s>2$. For the actual defining polynomial $F=x^2+y^2+z^2+xyz-4$, its actual partial derivative satisfies
\[
 F_x(2,-s,s)=4-s^2\ne0.
\]
Thus the family lies in the Jacobian-regular hypersurface locus, not just at singular nodes. Lean proves injections from $\mathbb N$ into both loci and uses genuine cardinalities to exclude \emph{every} finite count. It never uses the default value of \texttt{Nat.card} on an infinite type. After scalar extension to $\mathbb C$, actual trace-conjugation invariance proves that different family members are not simultaneously $\mathrm{SL}_2(\mathbb C)$-conjugate. The full scheme-smoothness criterion is not rederived.

\newpage
\section*{3. Exact finite-field count for all primes}
Let $K$ be any finite field of odd characteristic and $q=\#K$. With the same $+xyz$ convention, the invertible coordinate change $w=2z+xy$ gives
\[
 (x,y,z)\in S(K)\quad\Longleftrightarrow\quad
 w^2=(x^2-4)(y^2-4).
\]
Let $\chi$ be the actual quadratic character, extended by $\chi(0)=0$. Mathlib's square-root counting theorem gives $\#\{v:v^2=t\}=1+\chi(t)$. The actual hyperbola
\[
 v^2=x^2-4\quad\Longleftrightarrow\quad (x-v)(x+v)=4
\]
is in bijection with $K^\times$: take $h=x-v$; the inverse is
\[
 x=(h+4/h)/2,\qquad v=(4/h-h)/2.
\]
It has $q-1$ points. Counting the same actual fibers yields
$\sum_{x\in K}\chi(x^2-4)=-1$. Therefore, by multiplicativity,
\begin{align*}
 \#S(K)
 &=\sum_{x,y\in K}\bigl(1+\chi((x^2-4)(y^2-4))\bigr)\\
 &=q^2+\left(\sum_{x\in K}\chi(x^2-4)\right)^2=q^2+1.
\end{align*}
All bijections, finite fiber sums, cardinality casts and character identities are formalized; no desired counts are assigned to a named surface. Characteristic two is treated separately: Lean exhaustively counts the actual eight triples in $\mathbb F_2^3$, of which five satisfy the cubic. Hence $\#S(\mathbb F_p)=p^2+1$ for every prime $p$, including $2$. This counts the full affine character surface, including its singular points; it is not a separate assertion about the smooth finite-field locus.

\section*{4. Formalization, earlier submission and limits}
The Lean project uses actual \texttt{Matrix.SpecialLinearGroup}, \texttt{PresentedGroup}, \texttt{MvPolynomial.pderiv}, finite-field \texttt{quadraticChar}, explicit subtype equivalences, and \texttt{Cardinal.mk}. Its finite point-count interface requires a finite field carrier; it cannot silently turn the rational surface into a finite zero count. The final statements separate the rational impossibility from the positive finite-field theorem. All named local definitions and theorems are included in the axiom audit. The package supplies pinned dependencies, a new build, warning-as-error audit, complete fresh kernel replay, and source-bound receipts.

Earlier closed \href{https://github.com/The-Last-Math-Competition/The-Last-Math-Competition/pull/217}{PR 217} proposed finite-field cubic enumeration. The \href{https://github.com/The-Last-Math-Competition/The-Last-Math-Competition/pull/217#issuecomment-5965634432}{maintainer review} rejected the missing identification with the actual PVI monodromy object and specifically recognized the Cayley parameter $\kappa=4$ and its matching finite-field pattern. We attribute that exposed direction. This submission repairs the object bridge using the published PVI trace model and actual representations, proves the pattern for all primes rather than a sample, and separately closes the literal rational-point obstruction. We do not inherit unrelated material in that PR body or claim that the correct finite-field pattern is a counterexample.

The source's missing parameter choices and $\mathbb Q/\mathbb F_p$ mismatch are disclosed, not silently resolved. The result is a genuine rational counterexample to a universal finite rational-point reading, plus a complete proof of the anticipated finite-field formula for the standard Cayley member. It is not a classification of every PVI monodromy surface or a resolution of alternative unspecified arithmetic counting problems.

\textbf{AI disclosure.} Research, proof construction, Lean debugging and document preparation were led by a Codex worker, with separate AI supervision of source semantics, proof scope and PDF layout. No human verification or organizer acceptance is asserted. Public ideas and library results are attributed above.

\begin{thebibliography}{1}
\bibitem{CL} S.~Cantat and F.~Loray, \emph{Dynamics on Character Varieties and Malgrange irreducibility of Painlev\'e VI equation}, Ann. Inst. Fourier \textbf{59}(7) (2009), 2927--2978. Equations (1.3)--(1.5), (1.12), printed pp.~2928--2929, 2932. \href{https://doi.org/10.5802/aif.2512}{doi:10.5802/aif.2512}.
\end{thebibliography}
\end{document}
